% ch14 — Adaptive Mesh Refinement
\chapter{Adaptive Mesh Refinement}
\label{ch:amr}

Binary black-hole simulations span length scales from the near-horizon region
to the wave zone, requiring vastly different resolutions. This chapter develops
the Berger-Oliger adaptive mesh refinement (AMR) framework that places
fine grids only where needed.

\section{Fixed Mesh Hierarchy}
\label{sec:amr-hierarchy}
\coderef{amr}{rs}

We describe the nested box-in-box grid hierarchy with 2:1 refinement ratios,
define the level structure, and explain how grids are centred on each black
hole.

\section{Prolongation}
\label{sec:amr-prolongation}

We detail the interpolation (prolongation) operators that fill fine-grid ghost
zones from coarse-grid data, covering both polynomial and WENO options.

\section{Restriction}
\label{sec:amr-restriction}

We describe the restriction (averaging) operators that project fine-grid
solutions back onto the coarse grid to maintain consistency across levels.

\section{Subcycling in Time}
\label{sec:amr-subcycling}

We explain the recursive Berger-Oliger time-stepping algorithm that advances
fine levels with smaller time steps, synchronising at coarse-fine boundaries.

\section{Ghost Zones and Boundary Communication}
\label{sec:amr-ghost-zones}

We discuss the ghost-zone filling strategy, including inter-level and
intra-level communication patterns, and the special treatment required at
outer boundaries.

\paragraph{Key References.}
The AMR framework follows \cite{Berger1984,Berger1989}; the numerical relativity
application draws on \cite{Schnetter2004} and the Carpet driver
documentation; the prolongation operators follow \cite{Baumgarte2010}.
